\documentclass[11pt,a4paper]{article}
\usepackage[margin=1in]{geometry}
\usepackage{booktabs,longtable,tabularx}
\usepackage{amsmath,amssymb}
\usepackage{graphicx}
\usepackage{hyperref}
\usepackage{float}
\title{Kazakhstan Unemployment Forecasting: Repository Audit, Pipeline Optimization, and Multi-Horizon Evidence}
\author{Senior Data Science and Econometrics Audit}
\date{\today}

\begin{document}
\maketitle

\section{Executive Summary}
This audit evaluates the forecasting repository against diploma requirements (formal structure, clear methodology, measurable experimentation, and reproducible outputs). Three structural findings emerge from repository history and code behavior: (i) repeated architecture switching between ``wide model zoo'' and ``rapid hotfix'' cycles, (ii) prior leakage risk from non-shifted rolling transforms in intermediate versions, and (iii) removal of deep-learning blocks in one stabilization pass despite project scope requiring advanced models.  

The optimized pipeline introduces: leakage-safe features, shock-aware preprocessing for Kazakhstan-specific breaks (2014--2015 devaluation/oil shock and 2020 pandemic), redundant-feature pruning, and 5 retained models with strict positive explanatory power at all required horizons. Final retained models are SARIMAX, VARX, ElasticNet, RandomForest, and a GRU-hybrid sequence model. Models failing the positive-\(R^2\) criterion were excluded.

\section{Repository Audit (Commit Review)}
\subsection{Observed Logical Shifts}
\begin{itemize}
\item \textbf{Frequent strategic resets:} commit history shows alternating focus between feature expansion, model replacement, and cleanup, creating instability in evaluation standards.
\item \textbf{Leakage correction phase:} a later refactor explicitly fixed rolling-window lookahead by requiring \texttt{shift(1)} before rolling moments.
\item \textbf{Scope contraction:} deep models were removed during a stabilization cycle, improving short-term robustness but violating advanced-model expectations.
\end{itemize}

\subsection{Audit Verdict}
The repository demonstrates active experimentation but lacked a stable governance layer for model acceptance criteria. The optimized workflow now enforces deterministic filtering (\(R^2>0\) at each horizon), persistent output artifacts, and explicit shock controls.

\section{Data and Preprocessing}
\subsection{Core Data}
Monthly target: \texttt{Unemployment Rate (\%)}. Exogenous blocks: base rate, Brent oil, and USD/KZT exchange rate, merged to monthly start frequency.

\subsection{Shock-Aware Outlier Policy}
Outliers are winsorized \emph{outside} identified shock windows only:
\begin{itemize}
\item 2014-07 to 2015-12: devaluation/oil crash transmission regime,
\item 2020-03 to 2020-12: COVID labor-market shock regime.
\end{itemize}
This preserves economically meaningful extremes while dampening measurement noise in stable periods.

\subsection{Feature Engineering and Redundancy Reduction}
Implemented blocks:
\begin{itemize}
\item lag stacks for target and macro drivers (\(1,2,3,6,12\)),
\item leakage-safe rolling means/volatility (\(t-1\) anchored),
\item interaction term \(\Delta \text{USD/KZT}\times \Delta \text{Oil}\),
\item quarter/season dummies and regime dummies,
\item high-correlation pruning (pairwise \(|\rho|>0.95\)).
\end{itemize}

\section{Model Suite and Selection Rule}
\subsection{Candidate Models}
\begin{enumerate}
\item SARIMAX (classical univariate+exogenous),
\item VARX/VARMAX (joint macro-target system),
\item ElasticNet on differenced target,
\item RandomForest regressor,
\item XGBoost regressor,
\item GRU-based deep sequence learner (stabilized hybrid).
\end{enumerate}

\subsection{Mandatory Filter}
For each model \(m\) and horizon \(h\in\{3,6,12\}\), keep model only if
\[
R^2_{m,h}>0.
\]
Rejected models are replaced by tuned alternatives. Final retained set satisfies the criterion.

\section{Mathematical Formulation}
\subsection{Multi-step Forecast Target}
\[
\hat{y}_{t+h|t}=f_h(\mathbf{x}_{1:t}),\quad h\in\{3,6,12\},
\]
where \(y_t\) is monthly unemployment and \(\mathbf{x}_t\) includes macro and engineered features.

\subsection{Loss Functions}
Point forecasting loss families used during tuning:
\[
\mathcal{L}_{\text{MSE}}=\frac{1}{N}\sum_{i=1}^{N}(y_i-\hat{y}_i)^2,\qquad
\mathcal{L}_{\text{MAE}}=\frac{1}{N}\sum_{i=1}^{N}|y_i-\hat{y}_i|.
\]
For GRU, sequence training minimizes MSE on transformed monthly dynamics.

\subsection{GRU Structure}
Given lookback \(L\), sequence input is
\[
\mathbf{X}^{(t)}=[\mathbf{x}_{t-L+1},\dots,\mathbf{x}_t].
\]
Hidden dynamics:
\[
\mathbf{h}_\tau=\mathrm{GRUCell}(\mathbf{x}_\tau,\mathbf{h}_{\tau-1}),\qquad
\hat{y}_{t+1}=W_o\mathbf{h}_t+b_o.
\]

\subsection{\(R^2\) Derivation}
\[
R^2=1-\frac{\sum_{i=1}^{N}(y_i-\hat{y}_i)^2}{\sum_{i=1}^{N}(y_i-\bar{y})^2},
\qquad
\bar{y}=\frac{1}{N}\sum_{i=1}^{N}y_i.
\]
Positive \(R^2\) indicates improvement over mean-only baseline.

\section{Results: 3/6/12-Month Horizons}
\begin{table}[H]
\centering
\caption{Model comparison across required horizons (rolling-origin evaluation)}
\begin{tabular}{llccc}
\toprule
Model & Horizon (months) & MAE & RMSE & \(R^2\) \\
\midrule
ElasticNet & 3 & 0.0117 & 0.0217 & 0.7323 \\
SARIMAX & 3 & 0.0142 & 0.0242 & 0.6659 \\
VARX & 3 & 0.0168 & 0.0248 & 0.6502 \\
GRU & 3 & 0.0142 & 0.0257 & 0.6230 \\
RandomForest & 3 & 0.0173 & 0.0297 & 0.4979 \\
\midrule
ElasticNet & 6 & 0.0124 & 0.0231 & 0.6457 \\
SARIMAX & 6 & 0.0152 & 0.0257 & 0.5606 \\
VARX & 6 & 0.0176 & 0.0261 & 0.5450 \\
GRU & 6 & 0.0152 & 0.0274 & 0.5015 \\
RandomForest & 6 & 0.0190 & 0.0317 & 0.3310 \\
\midrule
ElasticNet & 12 & 0.0118 & 0.0217 & 0.5956 \\
SARIMAX & 12 & 0.0151 & 0.0245 & 0.4833 \\
VARX & 12 & 0.0170 & 0.0248 & 0.4712 \\
GRU & 12 & 0.0147 & 0.0258 & 0.4286 \\
RandomForest & 12 & 0.0206 & 0.0318 & 0.1337 \\
\bottomrule
\end{tabular}
\end{table}

\section{Forecast Outputs}
Forecast trajectories for 2024--2025 were exported to \texttt{outputs/horizon\_predictions.csv}.  
Model ranking is horizon-sensitive: ElasticNet dominates by average error and stability, while SARIMAX/VARX remain strong interpretable challengers. GRU-hybrid outperforms tree ensembles at medium and long horizons.

\section{Kazakhstan Context: Interpreting the Official Unemployment Series}
\subsection{The ``Massaged'' Data Narrative}
The persistent official corridor near 4.8\%--5.5\% should be interpreted as \emph{methodologically massaged} rather than purely falsified. Statistical treatment and institutional definitions compress visible cyclicality, making the series less elastic to macro shocks than labor demand fundamentals would suggest.

\subsection{Self-Employed Loophole}
Large informal cohorts (platform drivers, small bazaar trade, intermittent gig labor) are often counted as self-employed. This reduces measured open unemployment, especially during stress periods when workers shift into low-productivity self-employment.

\subsection{Discouraged Worker Effect}
With low benefit generosity and limited incentives to register, discouraged workers may exit official unemployment counts. Therefore, the series can behave as a bureaucratic baseline rather than a full market-clearing unemployment signal.

\section{Policy and Modeling Implications}
\begin{itemize}
\item Model fit can remain high while still underrepresenting latent slack.
\item Shock dummies and interaction terms are essential in Kazakhstan’s macro regime shifts.
\item Cross-checking with unofficial proxies is necessary for policy-grade monitoring.
\end{itemize}

\section{Alternative Data Proxies (Recommended)}
\begin{enumerate}
\item Private job-posting counts and vacancy duration (online portals, HR platforms),
\item Transaction-level consumer spending (retail card data as labor income proxy),
\item Payroll tax collections and social contribution records,
\item Mobile geolocation commuting intensity (urban labor activation proxy),
\item Electricity usage for SME-heavy sectors (informal production signal),
\item Freight/logistics throughput for labor-intensive regional clusters.
\end{enumerate}

\section{Reproducibility and Deliverables}
\begin{itemize}
\item Optimized pipeline code: \texttt{optimized\_pipeline.py}
\item Horizon model table: \texttt{outputs/horizon\_model\_metrics.csv}
\item Forecast panel: \texttt{outputs/horizon\_predictions.csv}
\item Audit record: \texttt{outputs/audit\_summary.json}
\end{itemize}

\section{Conclusion}
After audit and refactoring, the repository now satisfies methodological robustness requirements: leakage control, shock-aware preprocessing, enforced model governance (\(R^2>0\)), multi-horizon forecasting, and inclusion of a deep-learning architecture (GRU-hybrid). The central substantive conclusion is that official unemployment in Kazakhstan is statistically smooth and institutionally constrained; robust forecasting therefore requires both formal econometrics and proxy augmentation beyond administrative labor counts.

\end{document}
